\documentclass[11pt,a4paper]{article}

\usepackage[T1]{fontenc}
\usepackage[utf8]{inputenc}
\usepackage{lmodern}
\usepackage[margin=28mm]{geometry}
\usepackage{amsmath,amssymb,amsthm,mathtools}
\usepackage{booktabs}
\usepackage{microtype}
\usepackage[hidelinks]{hyperref}

\newcommand{\C}{\mathbb{C}}
\newcommand{\R}{\mathbb{R}}
\newcommand{\Z}{\mathbb{Z}}
\newcommand{\Chat}{\widehat{\mathbb{C}}}
\newcommand{\abs}[1]{\left|#1\right|}
\newcommand{\norm}[1]{\left\|#1\right\|}
\newcommand{\chord}{\chi}
\DeclareMathOperator{\tr}{tr}
\DeclareMathOperator{\Fix}{Fix}

\theoremstyle{plain}
\newtheorem{theorem}{Theorem}[section]
\newtheorem{proposition}[theorem]{Proposition}
\newtheorem{corollary}[theorem]{Corollary}
\theoremstyle{definition}
\newtheorem{definition}[theorem]{Definition}
\newtheorem{example}[theorem]{Example}
\theoremstyle{remark}
\newtheorem{remark}[theorem]{Remark}

\title{Dynamics of M\"obius Transformations\\
  \large Fixed points, conjugacy, and iteration on the Riemann sphere}
\author{}
\date{}

\begin{document}
\maketitle

\begin{abstract}
A M\"obius transformation is a rational map of degree one and a conformal
automorphism of the Riemann sphere. Despite its simple formula, its iterates
exhibit several distinct behaviors: periodic rotation, quasiperiodic rotation,
attraction and repulsion, and parabolic convergence. We derive the complete
classification from fixed points and conjugacy, relate it to matrix invariants,
and compute explicit examples. All convergence statements are made on the
sphere, so that escape to infinity is treated as convergence to an ordinary
point of the phase space.
\end{abstract}

\tableofcontents

\section{The phase space and the iteration problem}

The natural phase space is the Riemann sphere
\[
  \Chat=\C\cup\{\infty\}.
\]
A \emph{M\"obius transformation} is a map
\begin{equation}\label{eq:mobius}
  T(z)=\frac{az+b}{cz+d},
  \qquad a,b,c,d\in\C,\qquad \Delta=ad-bc\ne0.
\end{equation}
If $c\ne0$, its values at the pole and at infinity are
\[
  T(-d/c)=\infty,\qquad T(\infty)=a/c.
\]
If $c=0$, then $T(\infty)=\infty$. The determinant condition excludes
constant maps and ensures that $T$ is bijective, with inverse represented by
\[
  \begin{pmatrix}d&-b\\-c&a\end{pmatrix}.
\]
Thus both positive and negative iterates $T^n$, $n\in\Z$, are defined.
Here $T^0$ is the identity and $T^{n+1}=T\circ T^n$; the notation is
composition, not a pointwise power.

To discuss convergence at infinity, use the chordal metric
\[
  \chord(z,w)=\frac{\abs{z-w}}{\sqrt{1+\abs{z}^2}\sqrt{1+\abs{w}^2}},
  \qquad
  \chord(z,\infty)=\frac{1}{\sqrt{1+\abs{z}^2}}
\]
for finite $z,w$, with $\chord(\infty,\infty)=0$. In particular,
$z_n\to\infty$ on the sphere if and only if $\abs{z_n}\to\infty$.

The basic questions are: which points are fixed or periodic, which orbits
converge, and which geometric sets contain the orbits?

\section{Matrices, fixed points, and multipliers}

\subsection{Projective matrix representation}

Associate to \eqref{eq:mobius} the matrix
\[
  A=\begin{pmatrix}a&b\\c&d\end{pmatrix}\in\mathrm{GL}(2,\C).
\]
In homogeneous coordinates, $z=[z:1]$ and $\infty=[1:0]$, and
\[
  [u:v]\longmapsto[au+bv:cu+dv].
\]
Nonzero scalar multiples of $A$ represent the same transformation.
Composition corresponds to matrix multiplication, so $A^n$ represents $T^n$.
Consequently the group of M\"obius transformations is
$\mathrm{PGL}(2,\C)$.

\begin{proposition}\label{prop:fixed}
A nonidentity M\"obius transformation has either two distinct fixed points
or one fixed point of multiplicity two on $\Chat$.
\end{proposition}

\begin{proof}
A projective point is fixed precisely when its representing vector spans
an eigenline of $A$. A nonscalar complex $2\times2$ matrix has either two
distinct eigenlines or one eigenline and a nontrivial Jordan block.
Equivalently, the homogeneous fixed-point equation is
\[
  cu^2+(d-a)uv-bv^2=0.
\]
It has two roots on the projective line, counted with multiplicity, unless
it vanishes identically. The latter case means $c=b=0$ and $a=d$, which
is exactly the identity transformation.
\end{proof}

For finite fixed points the equation becomes
\begin{equation}\label{eq:fixed}
  cz^2+(d-a)z-b=0.
\end{equation}
When $c\ne0$, the two roots are
\[
  z_{\pm}=\frac{a-d\pm\sqrt{(a+d)^2-4\Delta}}{2c}.
\]
When $c=0$, infinity is fixed and must be included separately; merely
solving the finite equation can miss it.

\subsection{Multipliers in local coordinates}

\begin{definition}
The \emph{multiplier} at a fixed point $p$ is the derivative at zero of
$\phi\circ T\circ\phi^{-1}$, where $\phi$ is a holomorphic local
coordinate with $\phi(p)=0$.
\end{definition}

The chain rule shows that this definition does not depend on the chosen
coordinate and is invariant under M\"obius conjugacy. At a finite fixed
point $p$,
\begin{equation}\label{eq:derivative}
  T'(p)=\frac{\Delta}{(cp+d)^2}.
\end{equation}
At infinity, use $u=1/z$. For example, $T(z)=qz$ has multiplier $q$ at
zero and multiplier $q^{-1}$ at infinity, since
\[
  \frac{1}{T(1/u)}=\frac{u}{q}.
\]
A fixed point is called \emph{attracting} if the modulus of its multiplier
is less than one, \emph{repelling} if it is greater than one, and
\emph{neutral} if it equals one. Multipliers of M\"obius maps never vanish.

\section{Conjugacy and the two normal forms}

Two maps $T$ and $S$ are \emph{M\"obius conjugate} if
$S=h\circ T\circ h^{-1}$ for a M\"obius transformation $h$. Then
\begin{equation}\label{eq:conjugacy}
  T^n=h^{-1}\circ S^n\circ h,
\end{equation}
so $h$ transports fixed points, periodic points, and orbit limits.

\begin{theorem}[Normal forms]\label{thm:normal}
Every nonidentity M\"obius transformation is conjugate to exactly one of
the following types:
\begin{enumerate}
  \item $S(w)=qw$, where $q\in\C\setminus\{0,1\}$, if there are two
    distinct fixed points;
  \item $S(w)=w+1$, if there is one double fixed point.
\end{enumerate}
In the first case the parameter is determined up to replacement
$q\leftrightarrow q^{-1}$.
\end{theorem}

\begin{proof}
If the fixed points are $\alpha$ and $\beta$, choose $h$ sending them to
zero and infinity. When both are finite, one can take
\begin{equation}\label{eq:h}
  h(z)=\frac{z-\alpha}{z-\beta}.
\end{equation}
When $\beta=\infty$, take $h(z)=z-\alpha$; when $\alpha=\infty$,
take $h(z)=1/(z-\beta)$. The conjugated map fixes zero and infinity,
so it is $w\mapsto qw$. Its multiplier at zero is $q$, which equals
the multiplier of $T$ at $\alpha$. Interchanging the two fixed points
inverts $q$.

If there is only one fixed point $p$, first send $p$ to infinity.
The conjugated map fixes infinity and therefore has the affine form
$w\mapsto aw+b$. If $a\ne1$, it also fixes $b/(1-a)$, a contradiction.
Thus it is $w\mapsto w+b$, with $b\ne0$ because $T$ is not the identity.
Rescaling by $w\mapsto w/b$ gives $w\mapsto w+1$.
\end{proof}

For two distinct finite fixed points, \eqref{eq:conjugacy} gives an explicit
formula for every iterate:
\begin{equation}\label{eq:iterate}
  T^n(z)=\frac{\beta q^n h(z)-\alpha}{q^n h(z)-1}.
\end{equation}
The formula is interpreted on the sphere; a zero denominator represents
infinity, and values involving $h(z)=\infty$ are taken by continuity.
The reciprocal relation between the two fixed-point multipliers follows
immediately from the normal form.

\section{The dynamical classification}

We use \emph{loxodromic} for the entire class $\abs{q}\ne1$, and
\emph{hyperbolic} for its subclass with $q$ positive real. Some authors
reserve ``loxodromic'' for the nonhyperbolic members of this class.
This terminology convention does not change the dynamics.

\subsection{Loxodromic maps: attraction and repulsion}

Suppose $\abs{q}\ne1$. Interchange the fixed points if necessary so that
$0<\abs{q}<1$. In the coordinate $w=h(z)$,
\[
  S^n(w)=q^n w.
\]
Every finite $w$ tends to zero as $n\to+\infty$, while infinity remains
fixed. In backward time, every nonzero $w$ tends to infinity.

\begin{proposition}\label{prop:attraction}
A loxodromic map has an attracting fixed point $\alpha$ and a repelling
fixed point $\beta$. Its forward and backward basins are
\[
  \{z:T^n(z)\to\alpha\}=\Chat\setminus\{\beta\},\qquad
  \{z:T^{-n}(z)\to\beta\}=\Chat\setminus\{\alpha\}.
\]
Forward convergence is uniform on compact subsets of
$\Chat\setminus\{\beta\}$, and the analogous statement holds backward.
The only periodic points are the two fixed points.
\end{proposition}

\begin{proof}
In the normal coordinate, compact sets avoiding infinity are bounded, and
$\abs{q^n w}\le\abs{q}^n\sup\abs{w}\to0$ uniformly. Conjugation
transports this statement to the sphere. A nonzero finite periodic point
would satisfy $q^n w=w$ for some positive $n$, requiring $q^n=1$, which
is impossible when $\abs{q}\ne1$.
\end{proof}

Writing $q=\rho e^{i\theta}$, a nonzero orbit satisfies
\[
  \abs{w_n}=\rho^n\abs{w_0},\qquad
  \arg w_n\equiv\arg w_0+n\theta\pmod{2\pi}.
\]
It combines dilation with rotation. If $\theta\ne0$, the orbit lies on
the continuous logarithmic spiral
$w(t)=w_0\exp(t(\log\rho+i\theta))$, sampled at integer times.
For $q>0$, there is no rotation: each ray from zero is preserved. This
is the hyperbolic case. A negative real $q$ instead produces a half-turn
at every step and is not hyperbolic under our convention.

\subsection{Elliptic maps: periodic and irrational rotations}

Suppose $\abs{q}=1$ and $q\ne1$. Write $q=e^{i\theta}$. The map is
called \emph{elliptic}. Both fixed points are neutral, and every circle
$\abs{w}=r$, $0<r<\infty$, is invariant.

\begin{proposition}\label{prop:elliptic}
There are two possibilities for an elliptic map.
\begin{enumerate}
  \item If $\theta/(2\pi)=p/m\in\mathbb{Q}$ in lowest terms, then
    $T^m$ is the identity. Every point other than the two fixed points
    has exact period $m$.
  \item If $\theta/(2\pi)\notin\mathbb{Q}$, there are no nonfixed
    periodic points. Each other orbit is dense in the generalized
    circle $h^{-1}(\{\abs{w}=\abs{h(z)}\})$.
\end{enumerate}
\end{proposition}

\begin{proof}
A nonzero finite point is periodic exactly when $q^n=1$ for some $n>0$.
This yields the rational case and excludes additional periodic points in
the irrational case. For irrational $\theta/(2\pi)$, the subgroup
$\{e^{in\theta}:n\in\Z\}$ is dense in the unit circle: its closure is
an infinite closed subgroup of the circle, and every proper closed subgroup
of the circle is finite. Forward powers have the same closure. Multiplying
by $w_0$ and applying $h^{-1}$ proves the orbit-density claim.
\end{proof}

Here a \emph{generalized circle} means a Euclidean circle or a Euclidean
line together with infinity. M\"obius transformations send generalized
circles to generalized circles. Irrational elliptic orbits do not converge
to a fixed point; their entire invariant circle is their accumulation set.

\subsection{Parabolic maps: one neutral fixed point}

A nonidentity map with a double fixed point is called \emph{parabolic}.
In its normal coordinate,
\[
  S(w)=w+1,\qquad S^n(w)=w+n.
\]
Its unique fixed point is infinity, and its multiplier there is one:
\[
  u\longmapsto\frac{u}{1+u},\qquad
  \left.\frac{d}{du}\frac{u}{1+u}\right|_{u=0}=1.
\]

\begin{proposition}
If $T$ is parabolic with fixed point $p$, then
\[
  T^n(z)\longrightarrow p\quad\text{as }n\to+\infty
  \quad\text{and as }n\to-\infty
\]
for every $z\in\Chat$. The only periodic point is $p$. Convergence is
uniform on compact subsets of $\Chat\setminus\{p\}$, but not on
the whole sphere.
\end{proposition}

\begin{proof}
For finite $w$, $\abs{w+n}\to\infty$ in either direction, and this
convergence is uniform for bounded sets of $w$. No finite point satisfies
$w+n=w$ for a nonzero integer $n$. Uniform convergence on the entire sphere
is impossible: for each positive $n$, the initial point $w=-n$ maps to zero,
whose chordal distance from infinity is one.
\end{proof}

\begin{remark}[Convergence is not the same as attraction]
Although all parabolic orbits converge to $p$, the fixed point is neutral,
not attracting in the multiplier sense. It is not even Lyapunov stable:
for $S(w)=w+1$, initial points $w=-N$ can be arbitrarily close to infinity,
yet $S^N(-N)=0$. Thus pointwise convergence does not imply that nearby
orbits remain nearby at every intermediate time.
\end{remark}

Horizontal lines $\operatorname{Im}w=t$, completed by infinity, are
invariant under translation. Their inverse images under $h$ are invariant
generalized circles through the parabolic fixed point. Distinct members
of this family meet only at that point and are tangent there.

\subsection{The identity and a summary}

The identity must be separated from the parabolic class: it fixes every
point, and every orbit is constant. In particular, multiplier one alone
does not imply a nontrivial parabolic transformation.

\begin{center}
\small
\begin{tabular}{@{}lll@{}}
\toprule
Type & Normal form & Nonfixed forward orbits \\
\midrule
Identity & $w$ & None \\
Parabolic & $w+1$ & Converge to the double fixed point \\
Elliptic, rational angle & $e^{2\pi ip/m}w$ & Periodic of period $m$ \\
Elliptic, irrational angle & $e^{i\theta}w$ & Dense on invariant circles \\
Loxodromic & $qw$, $0<\abs{q}<1$ & Converge to the attracting point \\
\bottomrule
\end{tabular}
\end{center}
In the last row, the repelling fixed point itself remains fixed.

\section{Classification directly from matrix data}

The quantity
\begin{equation}\label{eq:tau}
  \tau(T)=\frac{(\tr A)^2}{\det A}
\end{equation}
is unchanged by rescaling $A$ and by matrix conjugation. If $A$ has
distinct eigenvalues $\lambda_1,\lambda_2$, the normal-form parameter is
their ratio in one order or the other, and
\begin{equation}\label{eq:traceq}
  \tau=q+q^{-1}+2.
\end{equation}
More specifically, the multiplier at the eigenline for $\lambda_1$ is
$\lambda_2/\lambda_1$: in an eigenvector basis the transverse coordinate
is multiplied by this ratio.

\begin{corollary}[Trace criterion]
For a nonidentity M\"obius transformation:
\begin{enumerate}
  \item $\tau=4$ if and only if the map is parabolic;
  \item $\tau\in[0,4)\subset\R$ if and only if the map is elliptic;
  \item $\tau\notin[0,4]\subset\R$ if and only if the map is loxodromic;
  \item $\tau\in(4,\infty)\subset\R$ if and only if the map is hyperbolic.
\end{enumerate}
The identity also has $\tau=4$, so it must be checked separately.
\end{corollary}

\begin{proof}
The repeated-eigenvalue condition is $(\tr A)^2=4\det A$.
A nonscalar matrix with this property gives the parabolic normal form.
For $q=e^{i\theta}\ne1$, \eqref{eq:traceq} gives
$\tau=2+2\cos\theta\in[0,4)$. Conversely, if $\tau\in[0,4)$,
the roots of $q^2-(\tau-2)q+1=0$ lie on the unit circle and are not one.
The remaining nonparabolic cases have $\abs{q}\ne1$.
Finally $q>0$, $q\ne1$, is equivalent to $q+q^{-1}>2$, or $\tau>4$.
\end{proof}

\begin{remark}
One may normalize $A$ to determinant one, in which case $\tau=(\tr A)^2$.
For complex matrices, however, the shortcut of comparing
$\abs{\tr A}$ with two does not classify the dynamics. For example,
$A=\operatorname{diag}(i/\sqrt{2},-i\sqrt{2})$ has determinant one and
$\abs{\tr A}=1/\sqrt{2}<2$, but its eigenvalue ratio is $-1/2$;
the transformation is loxodromic, not elliptic.
\end{remark}

\section{Worked examples}

\begin{example}[A hyperbolic map with finite fixed points]
Consider
\[
  T(z)=\frac{2z+1}{z+2}.
\]
Its fixed-point equation is $z^2-1=0$, so the fixed points are $1$ and
$-1$. Since $\Delta=3$,
\[
  T'(1)=\frac13,\qquad T'(-1)=3.
\]
For $h(z)=(z-1)/(z+1)$ a direct calculation gives
$h(T(z))=h(z)/3$. Consequently
\[
  T^n(z)=\frac{1+3^{-n}h(z)}{1-3^{-n}h(z)}.
\]
Every point except $-1$ converges to $1$ forward in time. The matrix
$\left(\begin{smallmatrix}2&1\\1&2\end{smallmatrix}\right)$ gives
$\tau=16/3>4$, confirming the hyperbolic classification.
\end{example}

\begin{example}[A genuinely rotating loxodromic map]
Let $T(z)=(i/2)z$. Zero is attracting, infinity is repelling, and
\[
  T^n(z)=2^{-n}i^n z.
\]
Each step rotates through $\pi/2$ while halving the radius. Every finite
orbit tends to zero. The iterates visit four rays in succession, showing
that the rotational part can be periodic even when the full orbit is not.
\end{example}

\begin{example}[Finite-order elliptic dynamics]
Let
\[
  T(z)=\frac{z-1}{z+1}.
\]
The fixed points are $i$ and $-i$. At $i$ the multiplier is
$2/(1+i)^2=-i$, a primitive fourth root of unity. Direct iteration gives
\[
  T^2(z)=-\frac1z,\qquad T^4(z)=z.
\]
Thus every nonfixed point has exact period four; for instance,
\[
  0\longmapsto-1\longmapsto\infty\longmapsto1\longmapsto0.
\]
Here $\tau=2$, in agreement with the elliptic trace criterion.
\end{example}

\begin{example}[An irrational elliptic orbit]
Let $T(z)=e^{2\pi i\sqrt{2}}z$. Zero and infinity are fixed, and every
other orbit is dense on the circle $\abs{z}=\abs{z_0}$. No other point
is periodic, because $n\sqrt{2}\notin\Z$ for every positive integer $n$.
The orbit stays bounded in the finite plane but does not converge.
\end{example}

\begin{example}[A parabolic map with a finite fixed point]
Consider
\[
  T(z)=\frac{z}{z+1}.
\]
Its only fixed point is zero, of multiplicity two. The coordinate
$h(z)=1/z$ gives $h(T(z))=h(z)+1$, hence
\begin{equation}\label{eq:parabolic-example}
  T^n(z)=\frac{z}{1+nz},\qquad n\in\Z.
\end{equation}
For fixed finite $z\ne0$, this tends to zero at the algebraic rate
$T^n(z)\sim1/n$, rather than the exponential rate of a loxodromic map.
Also $T^n(\infty)=1/n$ for $n\ne0$. A pole at an intermediate iterate
does not terminate the orbit: it merely sends that iterate to infinity.
Finally $T'(0)=1$ and $\tau=4$, so this is parabolic, not attracting in
the multiplier sense.
\end{example}

\section{A practical procedure and final perspective}

Given coefficients $a,b,c,d$, the following procedure determines the
dynamics without computing a long sequence of iterates.
\begin{enumerate}
  \item Verify $ad-bc\ne0$. If $b=c=0$ and $a=d$, the map is the identity.
  \item Solve the homogeneous fixed-point equation, remembering infinity.
  \item With two distinct fixed points, send them to zero and infinity;
    compute $q$ as the multiplier at the point sent to zero.
  \item If $\abs{q}\ne1$, identify the attracting and repelling points.
    If $\abs{q}=1$, decide whether $q$ is a root of unity.
  \item With one double fixed point, conjugate to translation.
  \item Use $T^n=h^{-1}\circ S^n\circ h$ for explicit orbits and limits.
    The invariant $\tau$ provides an independent classification check.
\end{enumerate}

The central observation is that degree-one dynamics are completely governed
by a dilation--rotation or a translation. There are no critical points,
no coexistence of distinct attracting cycles, and no nontrivial periodic
cycles except in the finite-order elliptic case. Nevertheless, the examples
separate several concepts that are easy to confuse: neutral and attracting
fixed points, bounded and convergent orbits, pointwise and uniform
convergence, and Euclidean escape versus convergence on the sphere.

\end{document}
